\documentclass[quiet]{ctexart}

\ExplSyntaxOn \makeatletter
\setCJKmainfont{Zhuque~Fangsong~(technical~preview)}[AutoFakeBold, ItalicFont = LXGW~WenKai]
\setCJKsansfont{LXGW~Marker~Gothic}[AutoFakeBold, AutoFakeSlant]
\setCJKmonofont{LXGW~WenKai}[AutoFakeSlant]
\ctexset{section/format+ = \raggedright}
\usepackage
  {
    geometry,mathtools,keytheorems,libertine,xeCJKfntef,
    siunitx,fixdif,derivative,physics2,tasks
  }
\usepackage[inline]{enumitem}
\usepackage{fontawesome5, fancyqr, tikz}
\hook_gput_code:nnn { cmd/faIcon/before } { . }
  {
    \group_begin:
    \fontsize { \fp_eval:n { .8 * \f@size } } { 0 } \selectfont
  }
\hook_gput_code:nnn { cmd/faIcon/after } { . } { \group_end: }
\geometry{margin = 1in}
\usepackage[warnings-off={mathtools-colon,mathtools-overbracket}]{unicode-math}
\setmathfont{STIX~Two~Math}
\usepackage[svgnames]{xcolor}
\newkeytheorem{note}
  [ style = remark, numbered = false,
    name = \color{DarkRed}\scshape\sffamily \faIcon{info-circle}~Note ]
\newkeytheorem{problem}
  [ style = definition,
    name = \color{DarkGreen}\scshape\sffamily \faIcon{pen-square}~Problem ]
\newkeytheorem{solution}
  [ style = definition, numbered = false,
    name = \color{MidnightBlue}\scshape\sffamily
           \faIcon{check-square}~Solution ]
\usephysicsmodule{ab}
\settasks { label = \sffamily\Alph*., item-indent = 2.5em, label-offset = .5pt }
\newcommand \true {\addtocounter{task}{1}\faIcon{check-circle}}
\def \@result#1 {\underline{~#1~}}
\def \s@result {\textsf{【\qquad 】}}
\DeclareRobustCommand \result {\@ifstar\s@result\@result}
\NewDocumentEnvironment{textfig}{O{7 \c_colon_str 2}}
  {
    \group_begin:
    \seq_set_split:Nee \l__whats_textfig_ratio_seq {\c_colon_str} {#1}
    \vspace{4.5pt plus 1.5pt minus 1.5pt} \noindent
    \minipage
      {
        \fp_eval:n
          {
            .96*\seq_item:Nn \l__whats_textfig_ratio_seq {1}/
              (
                \seq_item:Nn \l__whats_textfig_ratio_seq {1} + 
                \seq_item:Nn \l__whats_textfig_ratio_seq {2}
              )
          }\dimexpr \linewidth
      }
  }
  {
    \endminipage
    \vspace{4.5pt plus 1.5pt minus 1.5pt} \par
    \group_end:
  }
\NewDocumentCommand \textfigsplit {}
  {
    \endminipage
    \hspace{.02\linewidth}
    \minipage
      {
        \fp_eval:n
          {
            .96*\seq_item:Nn \l__whats_textfig_ratio_seq {2}/
              (
                \seq_item:Nn \l__whats_textfig_ratio_seq {1} + 
                \seq_item:Nn \l__whats_textfig_ratio_seq {2}
              )
          }\dimexpr \linewidth
      } 
  }

% \hook_gput_code:nnn { cmd/maketitle/before } { . } { \vskip -\baselineskip }
% \hook_gput_code:nnn { cmd/maketitle/after } { . }
%   { \__weyoung_maketitle_aftercode_aux: }
\cs_new_protected:Nn \__weyoung_maketitle_aftercode_aux:
  {
    \vspace*{-2\baselineskip}
    \begin{center}
      \begin{minipage}{.66\textwidth}
        \centering
        \begin{tasks}
          [ label-width = 5\ccwd, item-indent = 2cm,
            after-item-skip = 0pt ](2)
          \task[\textbf{课程信息}：] 理学院/A0715022
          \task[\textbf{授课教师}：] 陈许敏
          \task[\textbf{解析提供}：] \textsf{Skyrmion}
          \task[{\makebox[5em][s]{\textbf{\textsf\TeX~排版}：}}] \textsf{Skyrmion}
        \end{tasks}
      \end{minipage}
      \begin{minipage}{.15\textwidth}
        \centering
        \begin{tikzpicture}
          \node [ inner~sep = 0 ] (wechat) {
            \fancyqr
              [ size = 4em, image~padding = 1,
                image = \raisebox{4em/33}{\textcolor{DarkGreen}{\faIcon{weixin}}},
                r~color = DarkGreen, version = 4
              ]{http://weixin.qq.com/r/sRJiemTEGjllrSAZ90ej}
            };
          \node [ inner~sep = .5ex, align = center, right, font = \scriptsize ]
            at (wechat.east) {未\\央\\学\\社};
        \end{tikzpicture}
      \end{minipage}
      \begin{minipage}{.15\textwidth}
        \centering
        \begin{tikzpicture}
          \node [ inner~sep = 0 ] (qq) {
            \fancyqr
              [ size = 4em, image~x~padding = 2, image~y~padding = 1,
                image = \raisebox{4em/33}{\textcolor{MidnightBlue}{\faIcon{qq}}},
                l~color = MidnightBlue, version = 4
              ]{https://qm.qq.com/q/PH7CsQWvm2}
            };
          \node [ inner~sep = .5ex, align = center, left, font = \scriptsize ]
            at (qq.west) {物\\理\\营};
        \end{tikzpicture}
      \end{minipage}
    \end{center}
  }

\makeatother \ExplSyntaxOff
\usetikzlibrary{ fadings, patterns, arrows, decorations.markings }
\tikzset
  {
    > = stealth',
    every picture/.append style =
      { line join = round, line cap = round, thick }
  }

\title{\sffamily\bfseries 2023年HDU「大学物理A2」期末试题解答}
\author{} \date{}

\begin{document}

\maketitle

\section{选择题（每题3分，共36分）}

\begin{problem}
  一弹簧振子作简谐振动，其振动方程为 $x = A\cos(\omega t + \phi)$. 在求质点的振动动能时，
  得出下面4个表达式：
  \begin{tasks}[ label = \roman*. ](4)
    \task $\frac12 m\omega^2A^2 \sin^2(\omega t + \phi)$
    \task $\frac12 m\omega^2A^2 \cos^2(\omega t + \phi)$
    \task $\frac{2\pi^2}{T^2}mA^2 \sin^2(\omega t + \phi)$
    \task $\frac12kA^2 \cos^2(\omega t + \phi)$
  \end{tasks}
  其中 $m$ 是质点的质量，$k$ 是弹簧的劲度系数，$T$ 是振动的周期. 这些表达式中
  \begin{tasks}(4)
    \task i, iv 是对的.
    \task ii, iv 是对的.
    \task[\true] i, iii 是对的.
    \task ii, iii 是对的.
  \end{tasks}
\end{problem}
\begin{solution}
  质点振动动能
  $E_k = \frac12 m \dot x^2 = \frac12 m\omega^2A^2\sin^2(\omega t + \phi)$.
  由 $\omega = 2\pi/T$ 得 $E_k = \frac{2\pi^2}{T^2}mA^2\sin^2(\omega t + \phi)$.
\end{solution}

\begin{problem}
  小明同学利用一根长度为 $L$、劲度系数为 $k$ 的轻弹簧做简谐振动的研究试验，
  当弹簧下端挂一质量为 $m$ 的砝码时，测得系统周期为 $T$. 在实验过程中，
  小明不慎将弹簧一端打结发生变形，遂截去打结部分，剩余弹簧长度为 $\frac13L$.
  假如利用剩余部分继续做实验，要使得系统周期不变仍然为 $T$，则需要改挂质量为多少的砝码？
  \begin{tasks}(4)
    \task $\frac13\unit\m$
    \task[\true] $\qty3\m$
    \task $\qty9\m$
    \task $\frac19\unit\m$
  \end{tasks}
\end{problem}
\begin{solution}
  剩余弹簧的劲度系数为 $3k$，
  由弹簧振子周期 $T = 2\pi\sqrt{m/k}$ 得需改挂质量为 $3m$ 的砝码.
\end{solution}

\begin{textfig}[3:1]
  \begin{problem}
    一横波以波速 $u$ 沿 $x$ 轴负方向传播，$t$ 时刻波形图如图所示，则该时刻
    \begin{tasks}(2)
      \task A 点振动速度大于零.
      \task B 点静止不动.
      \task C 点向下运动.
      \task[\true] D 点振动速度小于零.
    \end{tasks}
  \end{problem}
  \textfigsplit \centering
  \begin{tikzpicture}[ scale = .8 ]
    \draw [ -> ] (-.9,0) -- (3.3,0) node [ below ] {$x$};
    \draw [ -> ] (0,-1.2) -- (0,1.2) node [ left ] {$y$};
    \draw (-.6,0) sin (0,.9) cos (.6,0) sin (1.2,-.9)
                  cos (1.8,0) sin (2.4,.9) cos (3,0);
    \shade [ ball color = black ] (.4,.45) circle (.06)
     node [ above right ] {$A$};
    \shade [ ball color = black ] (.6,0) circle (.06) node [ below left ] {$B$};
    \shade [ ball color = black ] (1.8,0) circle (.06)
     node [ below right ] {$C$};
    \shade [ ball color = black ] (2.8,.45) circle (.06)
     node [ above right ] {$D$};
    \draw [ <- ] (.8,1.2) -- node [ above ] {$u$} (1.6,1.2);
    \draw [ xshift = -.1cm, DarkGreen, dashed ]
      (-.6,0) sin (0,.9) cos (.6,0) sin (1.2,-.9)
              cos (1.8,0) sin (2.4,.9) cos (3,0);
  \end{tikzpicture}
\end{textfig}
\begin{solution}
  如图绿色虚线所示，下一刻 $A$ 点 $x$ 坐标对应的波形位置下移，$B$ 点 $x$ 坐标对应的波形位置下移，$C$ 点 $x$ 坐标对应的波形位置上移，$D$ 点 $x$ 坐标对应的波形位置下移.
\end{solution}

\begin{problem}
  光刻机中使用的光源都是需要进行色散矫正的浅紫外光和深紫外光.
  其中，浅紫外光的波长在 $\qty{365}\nm$ 至 $\qty{404}\nm$ 之间，
  深紫外光的波长在 $\qty{193}\nm$ 至 $\qty{248}\nm$ 之间.
  在光刻机中使用衍射光栅是为了对光做衍射和分光，将光分成不同频率的光线进行矫正，提高光学质量.
  若用衍射光栅对某一波长的单色光进行准确测定和分光，则在下列光栅中，选用哪一种最好
  \begin{tasks}
    \task 光栅单个狭缝宽度为 $\qty{0.5}\um$，狭缝间距为 $\qty1\mm$.
    \task 光栅单个狭缝宽度为 $\qty1\um$，狭缝间距为 $\qty1\mm$.
    \task[\true] 光栅单个狭缝宽度为 $\qty1\um$，狭缝间距为 $\qty2\mm$.
    \task 光栅单个狭缝宽度为 $\qty2\um$，狭缝间距为 $\qty4\mm$.
  \end{tasks}
\end{problem}
\begin{solution}
  由光栅的衍射方程 $d\sin\theta = k\lambda$ 可知，$d = \qty1\mm$ 和 $d = \qty4\mm$ 的光栅在观察紫外光时的角度过于密集在中心/分散在半径很大的外侧，均不利于观察. 所以选择光栅单个狭缝宽度为 $\qty2\mm$.
\end{solution}

\begin{textfig}[2:1]
  \begin{problem}
    小力同学在进行单缝夫琅禾费衍射实验时
    （装置如图所示，其中观察屏 $C$ 和光源 $S$ 始终放置于透镜焦平面位置），
    发现条纹太密不便观察，此时他应如何处理？
    \begin{tasks}
      \task 增大狭缝 $a$ 宽度
      \task 光源 $S$ 和狭缝之间选用焦距 $f$ 更大的透镜
      \task 光源 $S$ 和狭缝之间的透镜往狭缝方向移动
      \task[\true] 狭缝和观察屏 $C$ 之间选用焦距 $f$ 更大的透镜
    \end{tasks}
  \end{problem}
  \textfigsplit \centering
  \begin{tikzpicture}
    \filldraw [fill=gray] (0,1.5) rectangle (0.2,0.2);
    \filldraw [fill=gray] (0,-1.5) rectangle (0.2,-0.2);
    \draw [thick] (1.2,0)--(3,0);
    \draw [thick] (3,1.5)--(3,-1.5);
    \filldraw [fill=white] (1.5,0) ellipse (0.1 and 1.2);
    \foreach \y in {0.8,0,-0.8}
    \draw [thick,->,yshift=\y cm] (-1,0)--(-0.5,0);
    \draw [<->] (0,-1.6)--(0,-2)--(1.6,-2);
    \node [anchor=east] at (0,-1.6) {$y$};
    \node [anchor=north east] at (0,-2) {$O$};
    \node [anchor=west] at (1.6,-2) {$x$};
    \draw [ |<->| ] (1.5,-1.4) -- node [ above ] {$f$} (3,-1.4);
    \draw [<->] (0,0.2)--(0,-0.2) node [anchor=east,midway] {$a$};
    \node [anchor=south] at (-0.5,0.8) {$\lambda$};
    \node [anchor=south] at (1.5,1.2) {$L$};
    \node [anchor=west] at (3,1.5) {$C$};
\end{tikzpicture}
\end{textfig}
\begin{solution}
  根据单缝夫琅禾费衍射条纹间距公式 $\Delta x = 2f\lambda/a$
  可知其与入射光波长、狭缝宽度和观察屏前透镜焦距有关.
  所以狭缝和观察屏 $C$ 之间选用焦距 $f$ 更大的透镜以增大条纹间距，易于观察.
\end{solution}

\begin{problem}
  声纳技术是利用声波在水中的传播和反射特性，进行测距和导航的技术. 我国海军声纳技术，
  从早期仿制改进落后的苏式声纳系统到自行研制舰壳声纳和拖曳线阵列声纳，
  直至发展出主被动拖曳线阵列声纳，目前已处于世界领先行列.
  假设某潜艇静止在海底某处，潜艇声纳发送 $\qty{99}\kHz$ 的超声波，一不明物体以匀速朝潜艇运动，
  从该物体反射回来的超声波频率为 $\qty{101}\kHz$，
  忽略海水流速、温度、不均匀性等对超声波传播的影响，
  设海水中超声波的波速为 $u$，则该不明物体运动速度为
  \begin{tasks}(4)
    \task $\frac{1}{99}u$
    \task $\frac{1}{101}u$
    \task[\true] $\frac{1}{100}u$
    \task $99u$
  \end{tasks}
\end{problem}
\begin{note}
  提示：声波传播分为两个过程——潜艇发射声波时，不明物体为声波的接收者；不明物
  体反射声波时，该物体作为反射波的波源.
\end{note}
\begin{solution}
  已知多普勒效应观察者（Observer）和发射源（Source）的频率关系为
  $\nu = (u \pm v_o)\nu_0/(u \mp v_s)$.
  $v_o$ 为观察者速度，接近为 $+$，远离为 $-$；$v_s$ 为发射源速度，接近为 $-$，远离为 $+$.
  \begin{itemize}
    \item 不明物体接收到潜艇发射的声波频率为 $\nu_1 = (u + v)\nu_0/u$.
    \item 潜艇接到不明物体反射的声波频率为 $\nu_2 = u \nu_1/(u - v) = \qty{101}\kHz$.
    \item 带入 $\nu_0 = \qty{99}\kHz$ 得不明物体速度为 $v = u/100$.
  \end{itemize}
\end{solution}

\begin{problem}
  粒子在一维无限深方势阱中运动，波函数为
  $\psi(x) = \frac{1}{\sqrt a} \cos\frac{3\pi x}{2a},\ (-a \leq x \leq a)$.
  则粒子在 $x = 5a/6$ 出现的概率密度为
  \begin{tasks}(4)
      \task $1/a$
      \task[\true] $1/2a$
      \task $1/\sqrt{2a}$
      \task $1/\sqrt{a}$
  \end{tasks}
\end{problem}
\begin{solution}
  概率密度 $\omega(5a/6) = |\psi(5a/6)|^2 = 1/2a$. 
\end{solution}

\begin{problem}
  在不同的惯性系中，那一个物理量与参考系无关
  \begin{tasks}(4)
    \task 功
    \task 动能
    \task[\true] 力
    \task 质量
  \end{tasks}
\end{problem}

\begin{problem}
  常用激光笔发出的激光波长为 $\qty{650}\nm$，功率为 $\qty5\mW$，
  则激光笔每秒钟可发射的光子数约为
  \begin{tasks}(4)
    \task $\num{1.3e9}$
    \task $\num{4.9e24}$
    \task[\true] $\num{1.6e16}$
    \task $\num{1.2e39}$
  \end{tasks}
\end{problem}
\begin{solution}
  由光子能量 $E = nhc/\lambda$ 得 $n = \lambda E/hc = \num{1.636e16}$.
\end{solution}

\begin{problem}
  一束光是自然光和线偏振光的混合光，让它垂直通过一线偏振片. 若以此入射光束为轴旋转偏振片，
  测得透射光强度最大值是最小值的 $6$ 倍，那么入射光束中自然光与线偏振光的光强比值为
  \begin{tasks}(4)
    \task[\true] $2/5$
    \task $1/2$
    \task $2/7$
    \task $2/3$
  \end{tasks}
\end{problem}
\begin{solution}
  设自然光强度为 $I_1$，偏振光强度为 $I_2$.
  自然光透过偏振片后光强变为 $I_1/2$，偏振光透过偏振片后光强在 $0 \sim I_2$ 之间变化.
  所以投射光强最大值和最小值之比为 $(I_1/2 + I_2)/(I_1/2) = 6$，$I_1/I_2 = 2/5$.
\end{solution}

\begin{problem}
  北京正负电子对撞机是世界八大高能加速器中心之一，也是中国第一台高能加速器.
  高能加速器是利用电场和强磁场把带电粒子，如电子、质子加速到很高速度，然后去与靶物质相碰撞，
  碰撞的结果可产生大量的新的基本粒子，或新的现象.
  已知由高能加速器产生的一种不稳定粒子具有 $\qty{2.2}\us$ 的平均寿命，
  以 $0.6c$ 的速度穿过实验室. 则在该粒子发生衰变前，在实验室中飞行的平均距离为
  \begin{tasks}(4)
    \task $\qty{317}\m$
    \task[\true] $\qty{495}\m$
    \task $\qty{660}\m$
    \task $\qty{880}\m$
  \end{tasks}
\end{problem}
\begin{solution}
  在实验室参考系中，该粒子的平均寿命为 $\tau = \gamma\tau_0 = 1.25\tau_0$.
  所以其飞行的平均距离为
  $l = v\tau = 0.6c \times 1.25 \times \qty{2.2}\us = \qty{495}\m$.
\end{solution}

\begin{problem}
  静止质量不为零的微观粒子作高速运动，根据狭义相对论，粒子质量与粒子运动速度有关，
  则这时粒子德布罗意物质波的波长 $\lambda$ 与速度 $v$ 有如下关系
  \begin{tasks}(4)
    \task $\lambda \propto v$
    \task $\lambda \propto \frac1v$
    \task $\lambda \propto \sqrt{c^2 - v^2}$
    \task[\true] $\lambda \propto \sqrt{\frac1{v^2} - \frac1{c^2}}$
  \end{tasks}
\end{problem}
\begin{solution}
  由相对论能动量关系 $m^2c^4 = p^2c^2 + m_0^2c^4$ 得粒子动量
  $p = m_0c \sqrt{v^2/(c^2 - v^2)}$，所以波长
  $\lambda = h/p = \frac{h}{m_0} \sqrt{\frac1{v^2} - \frac1{c^2}}$.
  其中利用了 $m = \frac{m_0}{\sqrt{1 - v^2/c^2}}$.
\end{solution}

\section{填空题（每题3分，共18分）}

\begin{problem}
  两个同方向、同频率的简谐振动，其振动表达式分别为：$x_1 = 3\cos(\pi t + 5\pi/6)$ 和
  $x_2 = 4\cos(\pi t + \pi/3)$. 则合振动表达式为 \result{$5\cos(\pi t + 0.54\pi)$}.
\end{problem}
\begin{solution}
  很显然两振动相互垂直，所以合振动振幅 $A = \sqrt{A_1^2 + A_2^2} = 5$，
  合振动相位
  $ \phi = \pi + \arctan\ab(\frac{A_1\sin\phi_1 + A_2\sin\phi_2}
    {A_1\cos\phi_1 + A_2\cos\phi_2}) = 0.54\pi$. 所以合振动表达式为
  $y = 5\cos(\pi t + 0.54\pi)$.
\end{solution}

\begin{textfig}[4:1]
  \begin{problem}
    设沿弦线传播的入射波的表达式为
    $y_1 = A\cos\ab[2\pi(\frac tT - \frac x\lambda) + \phi]$，
    波在 $x = L$ 处（$B$ 点）发生反射，反射点为固定端（如图）.
    设波在传播和反射过程中振幅不变，则反射波的表达式为 $y_2 = $
    \result{$A\cos\ab[2\pi\ab(\frac tT + \frac {x - 2L}\lambda) + \phi + \pi]$}.
  \end{problem}
  \textfigsplit \centering
  \begin{tikzpicture}
    \draw [ -> ] (-.5,0) -- (2,0) node [ above ] {$x$};
    \draw [ -> ] (0,-.8) -- node [ below left ] {$O$} (0,.8)
     node [ right ] {$y$};
    \draw [ <-> ] (0,-.25) -- node [ below ] {$L$} (1.6,-.25);
    \fill [ pattern = north east lines, pattern color = gray ]
      (1.6,-.8) rectangle (1.8,.8);
    \draw [ line cap = butt ] (1.6,-.8) -- (1.6,.8);
    \fill (1.6,0) circle (.04) node [ above left ] {$B$};
  \end{tikzpicture}
\end{textfig}
\begin{solution}\leavevmode
  \begin{itemize}
    \item 将 $x = L$ 代入入射波表达式得 $B$ 点处的振动方程为
    $y_1(L,t) = A\cos\ab[2\pi\ab(\frac tT - \frac L\lambda) + \phi]$.
    \item 考虑半波损失，反射波在$B$ 点处的振动方程为
    $y_2'(L,t) = A\cos\ab[2\pi\ab(\frac tT - \frac L\lambda) + \phi + \pi]$.
    \item 加入传播项，反射波的表达式为
    $y_2(L,t) = A\cos\ab[2\pi\ab(\frac tT - \frac L\lambda + \frac {x'}\lambda) + \phi + \pi]$.
    \item 做坐标变换 $x' \to x - L$，以 $O$ 点为原点的反射波表达式为
    $y_2(L,t) = A\cos\ab[2\pi\ab(\frac tT + \frac {x - 2L}\lambda) + \phi + \pi]$.
  \end{itemize}
\end{solution}

\begin{problem}
  一等边三角形薄板相对地面沿其一条高线做高速运动，地面观察者观察到该板变为等腰直角三角形，
  则此时板相对地面运动速度为 \result{$\sqrt{2/3}c$}.
\end{problem}
\begin{solution}
  等边三角形在沿其运动方向发生尺缩，即 $\frac{l/2}{\sqrt3l/2} = \sqrt{1 - v^2/c^2}$，
  $v = \sqrt{2/3}c$.
\end{solution}

\begin{problem}
  一束自然光从折射率为 $1.4$ 的海水照射到钻石表面，当折射角为 $\qty{30}\degree$ 时，
  钻石表面的反射光为完全偏振光，则钻石的折射率为 \result{$2.42$}.
\end{problem}
\begin{solution}
  入射角 $\theta = \qty{90}\degree - \beta = \qty{60}\degree$.
  根据起偏角的定义 $\theta = \arctan(n_2/n_1)$ 得 $n_2 = n_1\tan\theta = 2.42$.
\end{solution}

\begin{problem}
  由两块长度均为 $\qty2\cm$ 的玻璃片（折射率 $n = 1.75$）所形成的空气劈形膜，
  其一端厚度为零，另一端厚度为 $\qty{0.002}\cm$. 现用波长为 $\qty{500}\nm$ 的单色平行光，
  垂直入射在膜的上表面，则形成的干涉暗条纹总数为 \result{$81$}.
\end{problem}
\begin{solution}
  由劈尖干涉公式，在玻璃片尾端 $2h = k\lambda$ 得尾端暗条纹级次 $k = 80$.
  考虑劈尖尖端的暗纹，形成的干涉暗条纹总数为 $81$.
\end{solution}

\begin{problem}
  在单缝衍射实验中，平行光垂直入射到狭缝，小明发现第 $3$ 级暗纹发生在衍射角 $\qty{30}\degree$
  的方向上，所用波长为 $\qty{600}\nm$，则小明所用单缝宽度为 \result{$\qty{3.6}\um$}
\end{problem}
\begin{solution}
  由单缝衍射公式 $a\sin\theta = k\lambda$
  得单缝宽度 $a = 3\lambda/\sin\theta = \qty{3.6}\um$.
\end{solution}

\section{计算题（共46分）}

\begin{problem}[本题5分]
  一质点按如下规律沿 $x$ 轴作简谐振动：$x = 0.2 \cos\ab(2\pi t + \frac13 \pi)$. 求此振动的周期、振幅、初相、速度最大值和加速度最大值.
\end{problem}
\begin{solution}
  将振动方程与 $y = A\cos(2\pi t/T + \phi)$ 比对，得周期 $T = \qty1\s$，
  振幅 $A = \qty{0.2}\m$，初相 $\phi = \pi/3$，
  速度最大值 $v_m = \omega A = 0.4\pi$，加速度最大值 $a_m = \omega^2 A = 0.8\pi^2$.
\end{solution}

\begin{problem}[本题8分]
  一平面简谐纵波沿着线圈弹簧传播. 设波沿着 $x$ 轴正向传播，
  弹簧中某圈的最大位移为 $\qty{1.0}\cm$，振动频率为 $\qty{10}\Hz$，
  弹簧中相邻两疏部中心的距离（即纵波波长）为 $\qty{12}\cm$. 当 $t = 0$ 时，
  在 $x = 0$ 处质元的位移为零并向 $y$ 轴负向运动. 试写出该波的表达式.
\end{problem}
\begin{solution}
  $t = 0$ 时 $x = 0$ 处质元的位移为零并向 $y$ 轴负向运动，
  所以振动初相位为 $\phi = \pi/2$. 波速 $u = \lambda f = \qty{1.2}{\m/\s}$.
  波的表达式为 $y = \num{1e-2} \cos\ab[20\pi\ab(t - x/1.2) + \pi/2]$.
\end{solution}

\begin{problem}[本题8分]
  小明在做双缝干涉实验的时候，想测一下常用激光笔的波长.
  于是，他用实验室波长 $\lambda = \qty{550}\nm$ 的单色平行光垂直入射到双缝上，
  在 $\qty2\m$ 远的屏幕上测量\emph{中央明纹两侧}的两条第 $5$ 级明纹中心的间距为 $\qty{0.11}\m$；
  保持装置不变，改用激光笔重复实验，
  测量中央明纹两侧的两条第 $5$ 级明纹中心的间距为 $\qty{0.13}\m$. 则
  \begin{enumerate}
    \item 该双缝干涉装置缝间距为多少？
    \item 所测激光笔波长为多少？
    \item 用厚度为 $e = \qty{5.5e-6}\m$、折射率为 $n = 1.3$ 的透明介质覆盖其中一条缝后，
    $\qty{550}\nm$ 光照下的零级明纹将移到原来的第几级明纹处？
  \end{enumerate}
\end{problem}
\begin{solution}\leavevmode
  \begin{enumerate}
    \item 由 $x_k = k\lambda D/d$ 得
    $d = 10 \lambda D/(x_5 - x_{-5}) = \qty{1e-4}\m$.
    \item 由 $x_k \propto \lambda$ 得激光笔波长
    $\lambda' = \lambda x_k'/x_k = \qty{650}\nm$.
    \item 用透明介质覆盖带来的光程差为 $\delta = (n - 1)e = \qty{1.65e-6}\m$，
    带来的极次移动 $\Delta k = \delta / \lambda = 3$.
    即$\qty{550}\nm$ 光照下的零级明纹将移到原来的第三级明纹处.
  \end{enumerate}
\end{solution}

\begin{problem}
  当氢原子从某初始状态跃迁到激发能（从基态到激发态所需的能量）为 $\Delta E = \qty{10.19}\eV$
  的状态时，发射的光为可见光. 利用每毫米有 $1200$ 条缝的光栅去观察该光（光线垂直入射到光栅上），
  在衍射角 $\qty{35.68}\degree$ 处观察到该光的第一亮纹，
  试求该氢原子初始状态的能量和主量子数，并回答处于该能级的氢原子最多可以发出几条谱线？
  （$\sin(\qty{35.68}\degree) = 0.5832$，$\qty1\eV = \qty{1.60e-19}\J$）
\end{problem}
\begin{solution}\leavevmode
  \begin{itemize}
    \item 光栅常数 $d = \qty1\mm/1200 = \qty{0.83}\um$.
    \item 由光栅方程 $d\sin\theta = k\lambda$，
    带入 $\theta = \qty{35.68}\degree$，$k = 1$
    得入射光波长 $\lambda = \qty{486.05}\nm$.
    \item 氢原子初始状态的能量
    $E_0 = -13.6 + \Delta E + hc/\lambda = -\qty{0.85}\eV$.
    主量子数为 $n = \sqrt{-13.6/E_0} = 4$.
    \item 处于该能级的氢原子最多可以发出 $C_4^2 = 6$ 条谱线.
  \end{itemize}
\end{solution}

\begin{problem}[本题6分]
  小力在做光电效应实验时，发现用波长为 $\lambda_1 = \qty{360}\nm$
  的紫光照射光阴极会产生光电流，假如给光阳极加上反向偏压则光电流将减小，
  直至电压增加到 $\qty{0.345}\V$ 时，光电流变为零. 则将该金属材料用作光电探测器探头时，
  其在可见光（$\qty{360}\nm$ 至$\qty{760}\nm$）范围内的工作波段是多少
  （即求可产生光电效应的波长范围）？
\end{problem}
\begin{solution}
  由光电效应方程 $eV = hc/\lambda - W$ 得金属的逸出功 $W = \qty{3.12}\eV$.
  则其产生光电效应的红限波长为 $\lambda = hc/W = \qty{397.58}\nm$，
  其在可见光范围内的工作波段为 $\qty{360}\nm \sim \qty{397.58}\nm$.
\end{solution}

\begin{problem}[本题5分]
  设光栅平面和透镜都与屏幕平行，在平面透射光栅上每条缝宽度为 $\qty{3.2}\um$，
  缝之间不透光区域为 $\qty{1.6}\um$，波长为 $\qty{600}\nm$ 的平行光垂直入射到光栅平面上，
  则求从中央主极大往外数（\emph{不包括中央主极大}），能观察到的第 3 条主极大条纹的衍射角为多少？
\end{problem}
\begin{solution}
  缺级 $k' = n(a + b)/a = 3,~ 6,~ 9,\, \ldots$. 所以能观察到的第 3 条主极大条纹的级次为 $4$.
  由光栅方程 $d\sin\theta = k\lambda$
  得 $\theta_4 = \arcsin[4\lambda/(a + b)] = \qty{30}\degree$.
\end{solution}

\end{document}